% titel: Blatt 3, A5: Ausbeute (Wasserstoff aus Zink)
% skript: Mortimer Kap. 4.5 (Ausbeute, Stöchiometrie)

% Teil I: Werkzeuge, kein Lösungsweg.

Bei der Ausbeute-Berechnung geht man von der eingesetzten Eduktmasse aus und
bestimmt zuerst die theoretisch mögliche Produktmasse, die dann mit der
tatsächlichen Ausbeute gewichtet wird.

\begin{werkzeugkasten}[Theoretische und tatsächliche Produktmasse]
Für die ausgeglichene Reaktion aus der Eduktmasse $m$ zuerst die Stoffmenge
des Edukts, daraus über die Stöchiometrie die \emph{theoretische} Stoffmenge
des Produkts und mit $m=n\cdot M$ dessen Masse:
\begin{equation}
  n_{\text{theor}}(\mathrm{H_2}) = n(\mathrm{Zn})\cdot
  \frac{\nu(\mathrm{H_2})}{\nu(\mathrm{Zn})}, \qquad
  m_{\text{theor}}(\mathrm{H_2}) = n_{\text{theor}}\cdot M(\mathrm{H_2}).
\end{equation}
Die tatsächlich erhaltene Masse ist
\begin{equation}
  m_{\text{real}} = \eta\cdot m_{\text{theor}},
  \qquad \eta = \qty{97.4}{\percent} = 0.974.
\end{equation}
\textit{Skript:} Mortimer Kap. 4 (Stöchiometrie, Ausbeute), Übung 3/5.
\end{werkzeugkasten}

\begin{achtung}[Darauf achten]
\begin{itemize}
  \item Ausbeute als Bruch ($0.974$) einsetzen, nicht als Prozent.
  \item Stöchiometrie prüfen: hier entsteht aus $1\ \mathrm{Zn}$ genau
        $1\ \mathrm{H_2}$ (Faktor 1:1).
  \item Erst die theoretische Masse (bei 100\%) berechnen, dann erst mit der
        Ausbeute multiplizieren; nicht den Faktor versehentlich vertauschen.
  \item Das Ergebnis (reale Masse) muss kleiner sein als die theoretische.
\end{itemize}
\end{achtung}
